\documentclass{article}%
\usepackage[T1]{fontenc}%
\usepackage[utf8]{inputenc}%
\usepackage{lmodern}%
\usepackage{textcomp}%
\usepackage{lastpage}%
\usepackage{geometry}%
\geometry{margin=2.2cm}%
%
\usepackage{amsmath}%
\usepackage{amssymb}%
\usepackage{booktabs}%
\usepackage{longtable}%
\usepackage{array}%
\usepackage{hyperref}%
\hypersetup{hidelinks}%
\title{IME 2026 {-} Matemática (claude\_sonnet\_5\_5) x Poliedro\\[4pt]\large Comparação com o gabarito Poliedro Resolve\thanks{Gerado com gpt-resolve: \url{https://github.com/lgabs/gpt-resolve}}}%
\date{10/11/2026}%
%
\begin{document}%
\normalsize%
\maketitle%
\section*{Resumo}\textbf{Nota média: 10,0/10} (10 de 10 itens com a resolução correta; cada item vale o mesmo dentro da questão).\par\medskip%
\begin{longtable}{@{}c c p{0.78\textwidth}@{}}\toprule Questão & Nota & Resumo \\ \midrule
1 & 10,0 & A resolução está correta e bate com o gabarito Poliedro: $\alpha$ = 3 e $\beta$ = 2, ou seja, z = 3 + 2i. A resolução usa o módulo ($\alpha$$^2$+$\beta$$^2$ = 13) e testa os sinais na expansão de z$^3$. O gabarito usa divisibilidade na parte real. Os dois caminhos são completos, provam que a solução é única e trazem a verificação. Não há divergência nem ponto a levar ao professor. \\
2 & 10,0 & A resolução está correta e bate com o gabarito do Poliedro: x = {-}7/2 ou x = {-}23/2. A transcrição confere com o enunciado do PDF original. Os dois usam o mesmo argumento: o que sobra sem (2x+7)\^{}3 é o binômio (x+7/2)\^{}4, e depois fatora{-}se. Conferi as contas com frações exatas. Não encontrei nenhuma divergência nem ponto duvidoso. \\
3 & 10,0 & Resolução e gabarito Poliedro concordam: k = (1 - a)(1 + $\surd$a), equivalente a (1 - $\surd$a)(1 + $\surd$a)$^2$. Conferi o enunciado no PDF original e refiz a conta analítica e numericamente, e o resultado está correto. A resolução é completa: justifica que o vértice está no domínio e relaciona a condição a $\neq$ 1/9 à exigência x $\neq$ a. O gabarito só tem deslizes de digitação ('k/4' no lugar de 4/k e '$\alpha$' no lugar de a), que não mudam a resposta. \\
4 & 10,0 & A resolução está correta e é equivalente ao gabarito do Poliedro. Os dois chegam às três retas pela origem, sem a origem: Im z/Re z = 1, -2+$\surd$3 ou -2-$\surd$3, isto é, arg z $\equiv$ 45$^\circ$, 105$^\circ$ ou 165$^\circ$ (mod 180$^\circ$). Só a parametrização é diferente: a resolução usa a parte imaginária como parâmetro, e o gabarito usa a parte real. Refiz a redução, que leva a (t-1)(t$^2$+4t+1) = 0, e conferi em Python que as duas famílias satisfazem a equação. A resolução ainda trata explicitamente as condições de existência (z não pode ser real nem imaginário puro) e faz uma verificação algébrica completa. Não há erro de nenhum dos lados. \\
5 & 10,0 & A resolução e o gabarito do Poliedro estão ambos corretos e são equivalentes. A resolução prova a desigualdade desenvolvendo o produto e aplicando x/y + y/x $\geq$ 2, com o lema demonstrado. É uma prova completa, que não depende de resultado externo e ainda identifica o caso de igualdade (a=b=c). O gabarito usa a MA $\geq$ MG e a demonstra numa observação. Conferi todas as identidades e a verificação numérica (49/4). O enunciado transcrito confere com o original. A questão está aprovada sem ressalvas que afetem a nota. \\
6 & 10,0 & A resolução e o gabarito do Poliedro dão o mesmo resultado, 6143850, e os dois estão certos. A sequência tem de ser uma PA de razão 3 (a\_j = 3j - 2), y\_j = 3 e x\_j+y\_j+z\_j = a\_j - 3 = 3j - 5, qualquer que seja a solução escolhida. Refiz a soma em Python e ela confere. A resolução é até mais rigorosa que o gabarito: prova que a PA é condição necessária e suficiente e mostra que a soma não depende da solução escolhida. O enunciado transcrito confere com o PDF original. \\
7 & 10,0 & A resolução está correta e bate com o gabarito do Poliedro: P(infectado | resposta negativa) = 0,045/0,144 = 5/16 = 31,25\%. Conferi com frações exatas. A resolução define bem os eventos, justifica os complementares e confere o resultado de duas formas: pelo evento complementar (11/16) e por uma contagem em 1000 computadores. A resolução do gabarito tem erros de digitação na notação (faltam barras de complemento e aparece 0,01$\cdot$0,99 onde deveria ser 0,1$\cdot$0,99), mas a resposta final 5/16 está certa. \\
8 & 10,0 & A resolução está correta e coincide com o gabarito Poliedro: a única solução real é (x, y, z) = (0, 0, 0). A resolução prova a unicidade pelo sinal de f(t) - t, com um argumento de cadeia que leva a absurdo nos casos x > 0 e x < 0. O gabarito usa a monotonia de f para chegar a f(x) = x. Os dois argumentos são rigorosos, e a resolução também confere a existência. Não há divergência. \\
9 & 10,0 & A resolução está certa e bate com o gabarito do Poliedro, só que escrita de outra forma. Ela dá r = a($\surd$3-1)(2-$\surd$2)/2 e o gabarito dá a($\surd$3-1)/(2+$\surd$2). As duas são iguais porque (2-$\surd$2)/2 = 1/(2+$\surd$2), e ambas valem $\approx$ 0,2144$\cdot$a. O raciocínio da resolução está completo: mostra por congruência que BM = DN, acha BM = a(2-$\surd$3) e descarta a raiz de fora, conclui que CMN é retângulo isósceles com catetos a($\surd$3-1) e calcula r = (b$_{1}$+b$_{2}$-h)/2. Não há erro em nenhum dos dois; o único defeito é uma passagem intermediária mal transcrita no gabarito, que não afeta o resultado. \\
10 & 10,0 & A resolução está correta e bate com o gabarito do Poliedro: V = 2V$_{0}$ = $\pi$$\surd$6$\cdot$a$^3$/108 $\approx$ 0,0713a$^3$. Os dois seguem o mesmo raciocínio: r = a$\surd$6/12, plano tangente passando pelos pontos médios das arestas, razão de semelhança 1/2, geração n com 4? esferas de volume V$_{0}$/8? e série geométrica de razão 1/2. Conferi as contas numericamente. A resolução ainda justifica que as esferas não se sobrepõem, o que o gabarito não faz. A única ressalva, de baixa importância, é a leitura do enunciado sobre incluir ou não a esfera inscrita original. \\
\bottomrule\end{longtable}%
\section*{Questão 1}%
\href{https://poliedroresolve.sistemapoliedro.com.br/vestibulares/ime/2026/ime-2-fase-dia-1-23-10-2025-2-fase-dia-1/dissertativa-1-matematica-2-fase-dia-1-ime-2026}{Resolução do gabarito}\par\smallskip%
\begin{longtable}{@{}p{0.06\textwidth} p{0.33\textwidth} p{0.33\textwidth} p{0.1\textwidth} p{0.08\textwidth}@{}}\toprule Item & Resolução & Gabarito & Veredito & Certo \\ \midrule
1 & $\alpha$ = 3 e $\beta$ = 2 (z = 3 + 2i) & $\alpha$ = 3 e $\beta$ = 2 (z = 3 + 2i) & igual & ambos \\
\bottomrule\end{longtable}%
\begin{itemize}%
\item%
\textbf{1)} As respostas são iguais. A transcrição bate com o enunciado original (pág. 1). Refiz as contas: |-9+46i|$^2$ = 81+2116 = 2197 = 13$^3$, então $\alpha$$^2$+$\beta$$^2$ = 13 e \{|$\alpha$|,|$\beta$|\} = \{2,3\}. Com |$\alpha$|=2 e |$\beta$|=3, a parte real dá -23$\alpha$ = $\mp$46, que é diferente de -9. Com |$\alpha$|=3 e |$\beta$|=2, sai -3$\alpha$ = -9, logo $\alpha$=3, e 23$\beta$ = 46, logo $\beta$=2. Confirmei que (3+2i)$^3$ = -9+46i. Uma busca exaustiva em Python com $\alpha$, $\beta$ $\in$ {[}-50,50{]} não nulos achou só (3,2). O caminho do gabarito é outro e também está certo: de $\alpha$$^3$ = 3($\alpha$$\beta$$^2$-3), $\alpha$ é múltiplo de 3 e divide -9, o que deixa $\alpha$ $\in$ \{$\pm$3, $\pm$9\}. Só $\alpha$=3 dá $\beta$ inteiro ($\beta$=$\pm$2), e a equação (II) fixa $\beta$=2.%
\end{itemize}%
\section*{Questão 2}%
\href{https://poliedroresolve.sistemapoliedro.com.br/vestibulares/ime/2026/ime-2-fase-dia-1-23-10-2025-2-fase-dia-1/dissertativa-2-matematica-2-fase-dia-1-ime-2026}{Resolução do gabarito}\par\smallskip%
\begin{longtable}{@{}p{0.06\textwidth} p{0.33\textwidth} p{0.33\textwidth} p{0.1\textwidth} p{0.08\textwidth}@{}}\toprule Item & Resolução & Gabarito & Veredito & Certo \\ \midrule
Questão 2 (soluções reais) & S = \{{-}23/2, {-}7/2\} (com {-}7/2 raiz tripla) & x = {-}7/2 ou x = {-}23/2 & igual & ambos \\
\bottomrule\end{longtable}%
\begin{itemize}%
\item%
\textbf{Questão 2 (soluções reais))} Os dois chegam às mesmas raízes pelo mesmo caminho. Passando {-}7\^{}4/16 para o primeiro membro, os termos sem (2x+7)\^{}3 formam (x+7/2)\^{}4, e a equação fica (x+7/2)\^{}4 + (2x+7)\^{}3 = 0. A resolução faz u = x+7/2 e obtém u\^{}3(u+8) = 0. O gabarito põe (2x+7)\^{}3 em evidência e obtém (2x+7)\^{}3{[}(2x+7)/16 + 1{]} = 0. Os dois métodos são equivalentes. Conferi com frações exatas: o polinômio original coincide com (x+7/2)\^{}4+(2x+7)\^{}3 em 5 pontos, logo são idênticos (grau 4). Ele se anula em x = {-}7/2 e em x = {-}23/2, e perto de {-}7/2 vale cerca de 8(x+7/2)\^{}3, o que confirma a raiz tripla. A observação sobre a multiplicidade vai além do que se pede, mas está correta.%
\end{itemize}%
\section*{Questão 3}%
\href{https://poliedroresolve.sistemapoliedro.com.br/vestibulares/ime/2026/ime-2-fase-dia-1-23-10-2025-2-fase-dia-1/dissertativa-3-matematica-2-fase-dia-1-ime-2026}{Resolução do gabarito}\par\smallskip%
\begin{longtable}{@{}p{0.06\textwidth} p{0.33\textwidth} p{0.33\textwidth} p{0.1\textwidth} p{0.08\textwidth}@{}}\toprule Item & Resolução & Gabarito & Veredito & Certo \\ \midrule
Questão 3 (valor de k em função de a) & k = (1 - $\surd$a)(1 + $\surd$a)$^2$ = (1 - a)(1 + $\surd$a) & k = (1 - a)(1 + $\surd$a) & igual & ambos \\
\bottomrule\end{longtable}%
\begin{itemize}%
\item%
\textbf{Questão 3 (valor de k em função de a))} As duas respostas coincidem; a resolução só acrescenta a forma fatorada equivalente, já que (1-$\surd$a)(1+$\surd$a)$^2$ = (1-a)(1+$\surd$a). Refiz a conta. Primeiro, f(x) = $\surd$x/(1-$\surd$x). Com s = $\surd$a e t = $\surd$x, a inclinação da secante é m = 1/{[}(1-s)(1-t)(t+s){]}. Para minimizar m, basta maximizar g(t) = (1-t)(t+s). O máximo é g = (1+s)$^2$/4, em t$_{0}$ = (1-s)/2 $\in$ (0, 1/2{]}. O ponto t$_{0}$ é admissível sempre que t$_{0}$ $\neq$ s, ou seja, a $\neq$ 1/9, que é justamente a restrição do enunciado. Assim, m\_min = 4/{[}(1-$\surd$a)(1+$\surd$a)$^2${]}. A reta ky - 4x = 2026 tem coeficiente angular 4/k, logo k = (1-a)(1+$\surd$a). Também conferi por varredura numérica em x (400 mil pontos) para a = 0; 0,04; 0,25; 0,3; 0,5; 0,81; 0,95. O mínimo bate com 4/k e ocorre em x = ((1-$\surd$a)/2)$^2$.%
\end{itemize}%
\paragraph{Pontos de dúvida}%
\begin{itemize}%
\item%
\textit{(key)} O texto do Poliedro diz que o coeficiente angular de ky - 4x = 2026 é 'k/4', mas o certo é 4/k (y = (4/k)x + 2026/k). Como a conta seguinte usa 4/k e chega ao k correto, é só um erro de digitação.%
\item%
\textit{(key)} Na resolução do Poliedro, o ponto aparece como '($\alpha$, f(a))' em vez de '(a, f(a))', e há frações sem parênteses na expressão da inclinação. É erro de digitação e não afeta o resultado.%
\item%
\textit{(statement)} O enunciado não justifica a exclusão a $\neq$ 1/9. A resolução explica o motivo: para a = 1/9 o mínimo cairia em x = a, que não define secante, e o ínfimo não seria atingido. O gabarito não comenta esse ponto. Uma banca rigorosa pode querer ver essa verificação, mas ela já está na resolução.%
\end{itemize}%
\section*{Questão 4}%
\href{https://poliedroresolve.sistemapoliedro.com.br/vestibulares/ime/2026/ime-2-fase-dia-1-23-10-2025-2-fase-dia-1/dissertativa-4-matematica-2-fase-dia-1-ime-2026}{Resolução do gabarito}\par\smallskip%
\begin{longtable}{@{}p{0.06\textwidth} p{0.33\textwidth} p{0.33\textwidth} p{0.1\textwidth} p{0.08\textwidth}@{}}\toprule Item & Resolução & Gabarito & Veredito & Certo \\ \midrule
Questão 4 (valores de z) & z = y(1+i), z = y($\surd$3-2+i) ou z = y(-2-$\surd$3+i), com y $\in$ $\mathbb{R}$*; ou seja, Re z/Im z = t $\in$ \{1, -2+$\surd$3, -2-$\surd$3\} & z = a(1+i) ou z = a{[}1+(-2$\pm$$\surd$3)i{]}, com a $\in$ $\mathbb{R}$*; ou seja, Im z/Re z $\in$ \{1, -2+$\surd$3, -2-$\surd$3\} & equivalente & ambos \\
\bottomrule\end{longtable}%
\begin{itemize}%
\item%
\textbf{Questão 4 (valores de z))} Os dois descrevem as mesmas três retas pela origem, sem a origem. Só muda o parâmetro: a resolução usa t = x/y (parte real sobre parte imaginária) e o gabarito usa b/a (parte imaginária sobre parte real). Como 1/(-2+$\surd$3) = -2-$\surd$3 e 1/(-2-$\surd$3) = -2+$\surd$3, os conjuntos coincidem. Em ângulos, as retas são $\theta$ = 45$^\circ$, 105$^\circ$ e 165$^\circ$, como a resolução diz na interpretação geométrica. Refiz a redução: com A = -it e B = -i(t$^2$-1)/(2t), a equação vira a+b+ab = 1, ou seja, t$^3$+3t$^2$-3t-1 = (t-1)(t$^2$+4t+1) = 0. É a mesma cúbica homogênea a$^3$+3a$^2$b-3ab$^2$-b$^3$ = 0 do gabarito. Também testei em Python as duas famílias com vários valores do parâmetro (positivos e negativos), e todas dão 1 no lado esquerdo. As condições de existência (Re z $\neq$ 0 e Im z $\neq$ 0) aparecem nos dois, de forma explícita na resolução e implícita no gabarito, por meio de a $\in$ $\mathbb{R}$* e de coeficientes angulares não nulos. A transcrição do enunciado confere com o PDF.%
\end{itemize}%
\paragraph{Pontos de dúvida}%
\begin{itemize}%
\item%
\textit{(key)} No gabarito, a letra 'a' é a parte real de z, e na resolução a letra 'y' é a parte imaginária. Por isso as duas formas finais parecem diferentes: ($\surd$3-2+i) contra (1+(-2-$\surd$3)i), por exemplo. Mas descrevem as mesmas retas. O corretor deve conferir pela razão Im z/Re z, não pela forma escrita.%
\item%
\textit{(key)} O gabarito escreve '-(a$^2$-b$^2$)/(2b$^2$i$^2$)' e, sem quebrar a linha, já passa a '+(a$^2$-b$^2$)/(2b$^2$)'. A conta está certa, porque i$^2$ = -1, mas a passagem está resumida e mal formatada. Além disso, o gabarito não declara explicitamente as condições de existência (a $\neq$ 0 e b $\neq$ 0), que só ficam garantidas por a $\in$ $\mathbb{R}$*.%
\end{itemize}%
\section*{Questão 5}%
\href{https://poliedroresolve.sistemapoliedro.com.br/vestibulares/ime/2026/ime-2-fase-dia-1-23-10-2025-2-fase-dia-1/dissertativa-5-matematica-2-fase-dia-1-ime-2026}{Resolução do gabarito}\par\smallskip%
\begin{longtable}{@{}p{0.06\textwidth} p{0.33\textwidth} p{0.33\textwidth} p{0.1\textwidth} p{0.08\textwidth}@{}}\toprule Item & Resolução & Gabarito & Veredito & Certo \\ \midrule
Questão 5 --- provar que (a+b+c)(1/a+1/b+1/c) $\geq$ 9 para a, b, c > 0 & Demonstrado. O produto se escreve como 3 + (a/b+b/a) + (a/c+c/a) + (b/c+c/b). Cada parcela é $\geq$ 2 pelo lema x/y + y/x - 2 = (x-y)$^2$/(xy) $\geq$ 0, que a própria resolução demonstra. Logo o produto é $\geq$ 9, com igualdade se e somente se a=b=c. Há também uma observação com uma segunda prova por MA $\geq$ MG. & Demonstrado por MA $\geq$ MG: a+b+c $\geq$ 3?(abc) e 1/a+1/b+1/c $\geq$ 3/?(abc). Como os dois lados são positivos, as desigualdades podem ser multiplicadas e dão $\geq$ 9 (c.q.d.). Como complemento, o gabarito demonstra a MA $\geq$ MG de três termos por p$^3$+q$^3$+r$^3$-3pqr = (p+q+r)$\cdot$$\frac12${[}(p-q)$^2$+(p-r)$^2$+(q-r)$^2${]} $\geq$ 0. & equivalente & ambos \\
\bottomrule\end{longtable}%
\begin{itemize}%
\item%
\textbf{Questão 5 --- provar que (a+b+c)(1/a+1/b+1/c) $\geq$ 9 para a, b, c > 0)} As duas provas chegam à mesma conclusão por caminhos diferentes. Ambas estão completas e corretas. A resolução desenvolve o produto e usa a soma de um número com seu inverso, x/y + y/x $\geq$ 2, por pares. É uma prova autocontida e mais elementar: não depende da MA $\geq$ MG sem demonstração. Ela ainda caracteriza corretamente a igualdade (a=b=c), o que o gabarito não faz, mas o enunciado não exige. Conferi as identidades com aritmética racional exata em 1000 trincas aleatórias: E - 9 = (a-b)$^2$/(ab) + (a-c)$^2$/(ac) + (b-c)$^2$/(bc), e p$^3$+q$^3$+r$^3$-3pqr = (p+q+r)$\cdot$$\frac12$$\Sigma$(p-q)$^2$. Conferi também a verificação numérica da resolução: 7$\cdot$(7/4) = 49/4 = 12,25. A prova principal do gabarito usa a MA $\geq$ MG como fato conhecido, mas a observação final a demonstra, então ele também fica completo. A transcrição do enunciado confere com o PDF original (página 1, 5ª questão).%
\end{itemize}%
\paragraph{Pontos de dúvida}%
\begin{itemize}%
\item%
\textit{(solution)} A 'Observação (outra prova)' da resolução usa a MA $\geq$ MG de três termos sem demonstrá{-}la. Isso não afeta a nota, porque a prova principal (Passos 1 a 3) é completa e autocontida. A observação é só um complemento.%
\item%
\textit{(key)} A linha principal do gabarito apoia{-}se na MA $\geq$ MG de três termos sem demonstração, e o próprio Poliedro alerta para uma possível penalização. A demonstração aparece na observação, e a fatoração está correta. Vale o professor confirmar se a banca do IME aceitaria a MA $\geq$ MG citada sem prova.%
\end{itemize}%
\section*{Questão 6}%
\href{https://poliedroresolve.sistemapoliedro.com.br/vestibulares/ime/2026/ime-2-fase-dia-1-23-10-2025-2-fase-dia-1/dissertativa-6-matematica-2-fase-dia-1-ime-2026}{Resolução do gabarito}\par\smallskip%
\begin{longtable}{@{}p{0.06\textwidth} p{0.33\textwidth} p{0.33\textwidth} p{0.1\textwidth} p{0.08\textwidth}@{}}\toprule Item & Resolução & Gabarito & Veredito & Certo \\ \midrule
Q6 --- soma de (x\_j+y\_j+z\_j) para j=1..2025 & 6143850 & 6 143 850 & igual & ambos \\
\bottomrule\end{longtable}%
\begin{itemize}%
\item%
\textbf{Q6 --- soma de (x\_j+y\_j+z\_j) para j=1..2025)} O enunciado transcrito confere com a página 2 do PDF. O sistema só é possível quando y = a\_\{j+1\} - a\_j = a\_\{j+2\} - a\_\{j+1\}. Valendo isso para j=1..2025, todas as 2026 diferenças consecutivas são iguais, então a sequência é uma PA. Como 10 = 1 + 3r, temos r = 3 e a\_j = 3j - 2. Pela 1ª equação, x+z = a\_j - 6, logo x+y+z = a\_j - 3 = 3j - 5. O gabarito escreve a\_j - 3 e a resolução escreve a\_\{j+1\} - 6, que é a mesma coisa porque a\_\{j+1\} = a\_j + 3. A soma de (3j-5) para j=1..2025 é 3$\cdot$2025$\cdot$1013 - 10125 = 6153975 - 10125 = 6143850. Pelo gabarito: (1+6073)$\cdot$2025/2 - 6075 = 6149925 - 6075 = 6143850. Conferi as duas contas em Python.%
\end{itemize}%
\paragraph{Pontos de dúvida}%
\begin{itemize}%
\item%
\textit{(key)} O gabarito diz que "uma solução possível" é a sequência ser uma PA, como se fosse só um caso entre vários. Na verdade a PA é obrigatória: subtraindo as equações vem a\_\{j+1\} - a\_j = a\_\{j+2\} - a\_\{j+1\} para todo j. É só uma imprecisão de redação, sem efeito no resultado. A resolução conferida prova a equivalência corretamente, nos dois sentidos.%
\item%
\textit{(statement)} Cada sistema tem infinitas soluções, porque só x+z fica determinado. Por isso o enunciado só faz sentido se x\_j+y\_j+z\_j não depender da solução escolhida. De fato não depende, já que x+y+z = a\_j - 3. A resolução justifica esse ponto explicitamente; o gabarito apenas o usa sem comentar.%
\end{itemize}%
\section*{Questão 7}%
\href{https://poliedroresolve.sistemapoliedro.com.br/vestibulares/ime/2026/ime-2-fase-dia-1-23-10-2025-2-fase-dia-1/dissertativa-7-matematica-2-fase-dia-1-ime-2026}{Resolução do gabarito}\par\smallskip%
\begin{longtable}{@{}p{0.06\textwidth} p{0.33\textwidth} p{0.33\textwidth} p{0.1\textwidth} p{0.08\textwidth}@{}}\toprule Item & Resolução & Gabarito & Veredito & Certo \\ \midrule
Q7 --- P(infectado | antivírus respondeu 'não infectado') & 5/16 = 0,3125 = 31,25\% & 5/16 & igual & ambos \\
\bottomrule\end{longtable}%
\begin{itemize}%
\item%
\textbf{Q7 --- P(infectado | antivírus respondeu 'não infectado'))} Os dois aplicam Bayes com P(I)=0,9, P(N|I)=1-0,95=0,05 e P(N|?)=1-0,01=0,99. Assim, P(I|N) = 0,045/(0,045+0,099) = 0,045/0,144 = 5/16. Conferi a conta em python com frações exatas e deu 5/16. A transcrição bate com o enunciado original da página 2.%
\end{itemize}%
\paragraph{Pontos de dúvida}%
\begin{itemize}%
\item%
\textit{(key)} Na resolução do gabarito faltam as barras de complemento: aparece 'P(B/A) = 0,01' quando deveria ser P(B|?), e o denominador está escrito como 'P(A)P(B/A) + P(A)P(B/A)'. É só erro de digitação ou de extração, porque a conta numérica (0,9$\cdot$0,05)/(0,9$\cdot$0,05 + 0,01$\cdot$0,99)\ldots{} na verdade usa 0,1$\cdot$0,99, e o resultado 5/16 está certo. Também está escrito '0,01 $\cdot$ 0,99' no lugar de '0,1 $\cdot$ 0,99', mas o valor final 5/16 só sai com 0,1$\cdot$0,99 = 0,099. Logo, isso também é erro de digitação e não muda a resposta.%
\item%
\textit{(statement)} A frase 'indica a presença de vírus quando o computador não está infectado com probabilidade de 1\%' foi lida como P(detecta | não infectado) = 0,01, que é a leitura padrão de taxa de falso positivo. A resolução e o gabarito adotam essa mesma leitura. Uma leitura alternativa, P(não infectado | detecta) = 0,01, seria pouco natural e mudaria o resultado.%
\end{itemize}%
\section*{Questão 8}%
\href{https://poliedroresolve.sistemapoliedro.com.br/vestibulares/ime/2026/ime-2-fase-dia-1-23-10-2025-2-fase-dia-1/dissertativa-8-matematica-2-fase-dia-1-ime-2026}{Resolução do gabarito}\par\smallskip%
\begin{longtable}{@{}p{0.06\textwidth} p{0.33\textwidth} p{0.33\textwidth} p{0.1\textwidth} p{0.08\textwidth}@{}}\toprule Item & Resolução & Gabarito & Veredito & Certo \\ \midrule
Q8 (unicidade e cálculo da solução) & Solução única (x, y, z) = (0, 0, 0). A unicidade é provada pelo sinal de f(t) - t = log(t + $\surd$(t$^2$+1)): é positivo para t > 0, negativo para t < 0 e nulo em t = 0. Se x > 0, a cadeia dá x > z > y > x; se x < 0, dá x < z < y < x. As duas são absurdas, logo x = 0 e daí y = z = 0. A existência é conferida por substituição direta. & Solução única (x, y, z) = (0, 0, 0). A prova usa que f é estritamente crescente: de x = f(f(f(x))) segue f(x) = x, logo x = y = z. Então log(x + $\surd$(x$^2$+1)) = 0, que leva a $\surd$(x$^2$+1) = 1 - x e a x = 0. & igual & ambos \\
\bottomrule\end{longtable}%
\begin{itemize}%
\item%
\textbf{Q8 (unicidade e cálculo da solução))} O enunciado confere com o original (pág. 2, 8ª questão). As respostas finais são idênticas. Os dois caminhos são válidos e diferentes. A resolução não precisa da monotonia de f: basta o sinal de f(t) - t, que ela prova com cuidado. Isso inclui o domínio, a desigualdade t + $\surd$(t$^2$+1) > 1 para t > 0 e o fato de g ser ímpar, porque ($\surd$(t$^2$+1) - t)($\surd$(t$^2$+1) + t) = 1. Assim ela exclui x > 0 e x < 0 e confere a existência. O gabarito usa a monotonia estrita de f para concluir f(x) = x. Se f(x) > x, então f(f(x)) > f(x) > x, e assim por diante. Depois ele resolve a equação de ponto fixo: x$^2$ + 1 = (1 - x)$^2$ dá x = 0, com 1 - x $\geq$ 0. Conferi numericamente: f(t) - t e f$^3$(t) - t têm o mesmo sinal de t, e o único zero é t = 0.%
\end{itemize}%
\paragraph{Pontos de dúvida}%
\begin{itemize}%
\item%
\textit{(key)} O gabarito afirma que f é estritamente crescente sem justificar. Isso é verdade: f'(t) = 1 + 1/$\surd$(t$^2$+1) > 0, ou f é soma de funções crescentes. Mas, numa prova discursiva de 'prove que', a falta dessa justificativa poderia custar algum desconto.%
\item%
\textit{(statement)} O enunciado não diz a base do log (decimal ou natural). Isso não muda nada, desde que a base seja maior que 1. A resolução registra isso explicitamente.%
\end{itemize}%
\section*{Questão 9}%
\href{https://poliedroresolve.sistemapoliedro.com.br/vestibulares/ime/2026/ime-2-fase-dia-1-23-10-2025-2-fase-dia-1/dissertativa-9-matematica-2-fase-dia-1-ime-2026}{Resolução do gabarito}\par\smallskip%
\begin{longtable}{@{}p{0.06\textwidth} p{0.33\textwidth} p{0.33\textwidth} p{0.1\textwidth} p{0.08\textwidth}@{}}\toprule Item & Resolução & Gabarito & Veredito & Certo \\ \midrule
Q9 --- raio do círculo inscrito no triângulo CMN & r = a($\surd$3 - 1)(2 - $\surd$2)/2 = a(2$\surd$3 - $\surd$6 + $\surd$2 - 2)/2 $\approx$ 0,2144$\cdot$a & r = a($\surd$3 - 1)/(2 + $\surd$2) & equivalente & ambos \\
\bottomrule\end{longtable}%
\begin{itemize}%
\item%
\textbf{Q9 --- raio do círculo inscrito no triângulo CMN)} As duas respostas são a mesma expressão escrita de outro jeito. Como (2 - $\surd$2)(2 + $\surd$2) = 2, vale (2 - $\surd$2)/2 = 1/(2 + $\surd$2). Assim, a($\surd$3 - 1)(2 - $\surd$2)/2 = a($\surd$3 - 1)/(2 + $\surd$2). Conferi com python3: as três formas (a fatorada da resolução, a expandida e a do gabarito) dão 0,214412717\ldots{} para a = 1, e área/p calculado direto também. O caminho é o mesmo nos dois. A resolução chega a BM = a(2 - $\surd$3) por Pitágoras (x$^2$ - 4ax + a$^2$ = 0) e o gabarito pela tg 15$^\circ$. Daí CM = CN = a($\surd$3 - 1), MN = CM$\cdot$$\surd$2, e r sai de r = (b$_{1}$ + b$_{2}$ - h)/2 na resolução e de {[}CMN{]}/p no gabarito. A resolução justifica BM = DN pelo caso cateto--hipotenusa e descarta a raiz a(2 + $\surd$3) porque cai fora do lado. A conferência numérica fecha: AM = MN $\approx$ 1,03528 para a = 1. O enunciado transcrito confere com o PDF original (página 2, 9ª questão).%
\end{itemize}%
\paragraph{Pontos de dúvida}%
\begin{itemize}%
\item%
\textit{(key)} Na última linha do gabarito, a fração do meio aparece como (a$^2$($\surd$3-1)$^2$)/(2+$\surd$2)a($\surd$3-1). Faltam os fatores 1/2 da área e do semiperímetro, que se cancelam, e os parênteses do denominador. O resultado final está certo; é só a passagem intermediária que ficou mal transcrita.%
\item%
\textit{(solution)} A resposta da resolução está em outra forma: fatorada a($\surd$3-1)(2-$\surd$2)/2 e expandida a(2$\surd$3-$\surd$6+$\surd$2-2)/2. O gabarito deixa o denominador sem racionalizar. Para um corretor ver de imediato que as duas coincidem, basta a identidade (2-$\surd$2)(2+$\surd$2) = 2.%
\end{itemize}%
\section*{Questão 10}%
\href{https://poliedroresolve.sistemapoliedro.com.br/vestibulares/ime/2026/ime-2-fase-dia-1-23-10-2025-2-fase-dia-1/dissertativa-10-matematica-2-fase-dia-1-ime-2026}{Resolução do gabarito}\par\smallskip%
\begin{longtable}{@{}p{0.06\textwidth} p{0.33\textwidth} p{0.33\textwidth} p{0.1\textwidth} p{0.08\textwidth}@{}}\toprule Item & Resolução & Gabarito & Veredito & Certo \\ \midrule
Questão 10 (volume total das esferas) & V = 2V$_{0}$ = ($\pi$$\surd$6/108)$\cdot$a$^3$ $\approx$ 0,0713 a$^3$, com V$_{0}$ = $\pi$$\surd$6$\cdot$a$^3$/216 & V = 2$\cdot$v$_{0}$ = ($\pi$$\surd$6/108)$\cdot$a$^3$, com v$_{0}$ = $\pi$$\surd$6$\cdot$a$^3$/216 & igual & ambos \\
\bottomrule\end{longtable}%
\begin{itemize}%
\item%
\textbf{Questão 10 (volume total das esferas))} Resolução e gabarito chegam ao mesmo resultado pelo mesmo caminho: h = a$\surd$6/3, r = h/4 = a$\surd$6/12 e V$_{0}$ = $\pi$$\surd$6a$^3$/216. O plano tangente fica a 2r = h/2 da base e corta as arestas nos pontos médios, então cada novo tetraedro é semelhante ao anterior com razão 1/2. Na geração n há 4? esferas, cada uma com volume V$_{0}$/8?. A soma é a série geométrica de razão 1/2, que dá 2V$_{0}$. Conferi em Python: V$_{0}$ $\approx$ 0,035626a$^3$, 2V$_{0}$ $\approx$ 0,071253a$^3$ = $\pi$$\surd$6/108$\cdot$a$^3$, cerca de 60,5\% do volume do tetraedro ($\surd$2/12$\cdot$a$^3$ $\approx$ 0,1179a$^3$). A resolução também justifica, por coordenadas baricêntricas, que as esferas não se sobrepõem, e o gabarito não faz isso. Há um modo mais direto de ver: o raio da esfera inscrita no octaedro central (aresta a/2), (a/2)$\cdot$$\surd$6/6 = a$\surd$6/12, é igual a r. Logo a esfera inscrita em ABCD é exatamente a esfera inscrita no octaedro e não entra nos tetraedros dos cantos.%
\end{itemize}%
\paragraph{Pontos de dúvida}%
\begin{itemize}%
\item%
\textit{(statement)} A expressão "todas as esferas interiores ao tetraedro ABCD" poderia ser lida como se excluísse a esfera inscrita original, já que ela é tangente às faces. Nessa leitura a resposta seria V$_{0}$ = $\pi$$\surd$6a$^3$/216. A resolução e o Poliedro incluem a esfera original e chegam a 2V$_{0}$, que é a leitura mais natural. Ainda assim, convém saber como o IME pontua quem respondeu só a soma das esferas novas.%
\item%
\textit{(key)} O gabarito não justifica que as esferas de gerações e de cantos diferentes têm interiores disjuntos, condição necessária para somar os volumes. A resolução cobre esse ponto. Isso não muda a resposta, mas pode pesar na correção de uma prova discursiva.%
\end{itemize}%
\end{document}